\chapter{Estimation}\label{ch:qm:estimation}

A researcher presents a strategy: over five years its daily P\&L averaged 4.0 basis points with a
\omterm{def:qm:estimation:estimator}{standard error} of 1.2, a $t$-statistic of 3.4, comfortably significant. A colleague notices that
each day's position is held for five days, so consecutive days share four fifths of their
positions, and recomputes the \omterm{def:qm:estimation:estimator}{standard error} allowing for the autocorrelation that creates: 2.3.
The $t$-statistic falls to 1.7, and the strategy is no longer distinguishable from luck. Nothing
was wrong with the average; the formula for its uncertainty assumed independent days. This chapter
is about \omterm{def:qm:estimation:estimator}{estimators} and their uncertainty: maximum likelihood and its asymptotics, what happens
when the model is wrong, \omterm{def:qm:estimation:m}{M-estimators} and moments, the \omterm{met:qm:estimation:delta}{delta method} that gives the \omterm{def:qm:estimation:sharpe}{Sharpe ratio} its
\omterm{def:qm:estimation:estimator}{standard error}, and the heteroskedasticity- and autocorrelation-consistent variances that repaired
the researcher's $t$-statistic.

\section{Maximum likelihood}

\begin{definition}[Estimator, unbiased, consistent, standard error]\label{def:qm:estimation:estimator}
An \emph{estimator}\index{estimator} $\hat\theta_n$ of a parameter $\theta_0$ is a function of the data $X_1, \dots, X_n$. It is an
\emph{unbiased estimator}\index{unbiased estimator} if $\E[\hat\theta_n] = \theta_0$, and a \emph{consistent
estimator}\index{consistent estimator} if $\hat\theta_n \xrightarrow{\P} \theta_0$ as $n \to \infty$. Its \emph{standard
error}\index{standard error} $\mathrm{se}(\hat\theta_n)$ is (an estimate of) its standard deviation.
\end{definition}

The sample variance with divisor $n$ is biased and consistent; with divisor $n - 1$ it is unbiased. Consistency
is what matters with the sample sizes of finance, and so does the \omterm{def:qm:estimation:estimator}{standard error}: an estimate without one is an
anecdote.

\begin{definition}[Likelihood function, maximum likelihood estimator, score, Fisher information]\label{def:qm:estimation:mle}
For independent observations with density $f(x; \theta)$, the \emph{likelihood function}\index{likelihood function} is
$\theta \mapsto \prod_if(X_i; \theta)$ and the log-likelihood $\ell_n(\theta) = \sum_i\ln f(X_i; \theta)$. The \emph{maximum
likelihood estimator}\index{maximum likelihood estimator} (MLE) maximises it. The \emph{score function}\index{score
function} is $\nabla_\theta\ln f(X; \theta)$, and the \emph{Fisher information}\index{Fisher information} is $\mathcal I(\theta) =
\E[\nabla\ln f\,\nabla\ln f^\top] = -\E[\nabla^2\ln f]$ per observation.
\end{definition}

\begin{theorem}[Asymptotics of maximum likelihood]\label{thm:qm:estimation:mle}
Under regularity conditions (identifiability, smoothness, a true value in the interior), the MLE is consistent and
\[
\sqrt n(\hat\theta_n - \theta_0) \xrightarrow{d} \mathcal N\bigl(0, \mathcal I(\theta_0)^{-1}\bigr).
\]
No \omterm{def:qm:estimation:estimator}{unbiased estimator} has a smaller variance than $\mathcal I(\theta_0)^{-1}/n$ (the Cramér--Rao bound), so the MLE is
asymptotically efficient.
\end{theorem}

\begin{proof}[Partial proof]
The score at $\theta_0$ has mean zero and variance $\mathcal I$, so $n^{-1/2}\nabla\ell_n(\theta_0) \to \mathcal N(0, \mathcal I)$ by the central
limit theorem. Expanding $0 = \nabla\ell_n(\hat\theta_n) = \nabla\ell_n(\theta_0) + \nabla^2\ell_n(\tilde\theta)(\hat\theta_n - \theta_0)$ and using
$-n^{-1}\nabla^2\ell_n \to \mathcal I$ gives $\sqrt n(\hat\theta_n - \theta_0) = \mathcal I^{-1}n^{-1/2}\nabla\ell_n(\theta_0) + o_{\P}(1)$. Consistency
and the Cramér--Rao bound are in Casella and Berger (2002).
\end{proof}

\begin{example}[The rate of arrivals]\label{ex:qm:estimation:rate}
For $n$ exponential waiting times with rate $\lambda$, $\ell_n(\lambda) = n\ln\lambda - \lambda\sum t_i$, so $\hat\lambda = n/\sum t_i$, $\mathcal I(\lambda) =
1/\lambda^2$, and $\mathrm{se}(\hat\lambda) = \hat\lambda/\sqrt n$. A hundred gaps summing to 50 seconds give $2 \pm 0.2$ arrivals per second.
The Hawkes fits of chapter~7 are the same computation with a harder likelihood.
\end{example}

\section{Misspecification and the sandwich}

Models are approximations; the question is what maximum likelihood estimates when the density is wrong.

\begin{definition}[Kullback--Leibler divergence, quasi-maximum likelihood]\label{def:qm:estimation:kl}
The \emph{Kullback--Leibler divergence}\index{Kullback--Leibler divergence} of a density $f$ from the true density $p$
is $D_{\mathrm{KL}}(p\,\|\,f) = \E_p[\ln(p(X)/f(X))] \ge 0$, zero only if $f = p$. \emph{Quasi-maximum
likelihood}\index{quasi-maximum likelihood} maximises a likelihood that is not believed to be the true one.
\end{definition}

\begin{proposition}[What a wrong likelihood estimates]\label{prop:qm:estimation:qmle}
Under regularity conditions the quasi-MLE converges to the pseudo-true value $\theta^*$ that minimises $D_{\mathrm{KL}}(p\,\|\,f_\theta)$,
and $\sqrt n(\hat\theta_n - \theta^*) \xrightarrow{d} \mathcal N(0, A^{-1}BA^{-1})$ with $A = -\E[\nabla^2\ln f_{\theta^*}]$ and $B = \Var(\nabla\ln
f_{\theta^*})$.
\end{proposition}

\begin{proof}
$n^{-1}\ell_n(\theta) \to \E_p[\ln f_\theta] = \E_p[\ln p] - D_{\mathrm{KL}}(p\,\|\,f_\theta)$, maximised at $\theta^*$. The expansion of the
previous proof holds with $\mathcal I$ replaced by $A$ in the Hessian and by $B$ in the score's variance; they coincide only when
the model is right (the information equality).
\end{proof}

\begin{definition}[Sandwich variance]\label{def:qm:estimation:sandwich}
The \emph{sandwich variance}\index{sandwich variance} of an \omterm{def:qm:estimation:estimator}{estimator} defined by an estimating equation is $A^{-1}BA^{-1}$:
the bread $A$, the expected derivative of the estimating equation, around the meat $B$, the variance of the equation
itself. Estimated by the sample Hessian and the outer product of the per-observation scores, it is the robust
(Huber--White) variance.
\end{definition}

Fit a normal model to 2\,000 draws of a Student $t$ with 6 degrees of freedom. The estimate of the log standard
deviation is still consistent for the log standard deviation, but the Hessian \omterm{def:qm:estimation:estimator}{standard error}, 0.0158, assumes normal
tails; the sandwich gives 0.0230, and over 4\,000 simulated samples the actual spread is 0.0237
(\cref{fig:qm:estimation:sandwich}). With fatter tails still (4 degrees of freedom, infinite fourth moment) the
variance of the sample variance does not exist, and no \omterm{def:qm:estimation:estimator}{standard error} is right.

\begin{omfigure}
\begin{tikzpicture}
\begin{axis}[width=11cm,height=4.8cm,
  xlabel={$\ln\hat\sigma - $ its mean},ylabel={density},
  tick label style={font=\small},label style={font=\small},
  xmin=-0.1,xmax=0.1,ymin=0,ymax=27,grid=major,grid style={gray!25},
  xtick={-0.1,-0.05,0,0.05,0.1},xticklabel style={/pgf/number format/fixed},scaled x ticks=false,
  legend columns=3,legend style={font=\footnotesize,at={(0.5,-0.3)},anchor=north,draw=none,fill=none,/tikz/every even column/.append style={column sep=8pt}},
  table/col sep=comma]
\addplot[black!60,thick,const plot mark mid] table[x=x,y=hist]{figdata/methods/11-estimation/sandwich.csv};
\addplot[omThm,very thick,dashed] table[x=x,y=hessian]{figdata/methods/11-estimation/sandwich.csv};
\addplot[omDef,very thick] table[x=x,y=sandwich]{figdata/methods/11-estimation/sandwich.csv};
\legend{4\,000 samples,Hessian (normal model),sandwich}
\end{axis}
\end{tikzpicture}

\omcaption{The log standard deviation estimated by a normal likelihood from 2\,000 Student-$t_6$ observations: its actual
sampling distribution (steps) against the normal approximations with the Hessian \omterm{def:qm:estimation:estimator}{standard error} (dashed), too narrow,
and with the sandwich (solid). Data: the chapter's tutorial, seeded.}\label{fig:qm:estimation:sandwich}
\end{omfigure}

\section{M-estimators and the method of moments}

\begin{definition}[M-estimator, method of moments, generalised method of moments]\label{def:qm:estimation:m}
An \emph{M-estimator}\index{M-estimator} maximises (or zeroes the derivative of) a sample average $n^{-1}\sum_im(X_i; \theta)$;
maximum likelihood takes $m = \ln f$, least squares $m = -(y - x^\top\theta)^2$. The \emph{method of
moments}\index{method of moments} solves $n^{-1}\sum_ig(X_i) = \E_\theta[g(X)]$ for as many moments as parameters. The
\emph{generalised method of moments}\index{generalised method of moments} (GMM) minimises $\bar g_n(\theta)^\top W\bar g_n(\theta)$ for
$\bar g_n(\theta) = n^{-1}\sum_ig(X_i; \theta)$ with more moment conditions than parameters and a weight matrix $W$.
\end{definition}

All three have the sandwich asymptotics of \cref{prop:qm:estimation:qmle}, with $A$ the derivative of the estimating
equations and $B$ their long-run variance; GMM is efficient when $W$ is the inverse of that variance (Hansen, 1982),
and the minimised objective then tests the overidentifying conditions. For an \omterm{def:qm:stochastic-differential-equations:ou}{Ornstein--Uhlenbeck process} sampled
daily, the \omterm{def:qm:estimation:m}{method of moments} on the first autocorrelation, $\hat\kappa = -\ln\hat\rho(1)/\Delta t$, coincides with the conditional
maximum likelihood of the autoregression (chapter~17), and both inherit its small-sample bias.

\section{The delta method and the Sharpe ratio}

\begin{method}[Delta method]\label{met:qm:estimation:delta}
If $\sqrt n(\hat\theta_n - \theta_0) \xrightarrow{d} \mathcal N(0, V)$ and $g$ is differentiable at $\theta_0$, then $\sqrt n(g(\hat\theta_n) - g(\theta_0))
\xrightarrow{d} \mathcal N(0, \nabla g^\top V\nabla g)$. The \emph{delta method}\index{delta method} linearises a smooth function of
an estimate and propagates its variance.
\end{method}

\begin{definition}[Sharpe ratio]\label{def:qm:estimation:sharpe}
The \emph{Sharpe ratio}\index{Sharpe ratio} of a strategy is its expected excess return per unit of standard deviation,
$\mathrm{SR} = \mu/\sigma$, per period; annualised by $\sqrt{\text{periods a year}}$ when returns are independent across periods. Its
estimate is $\widehat{\mathrm{SR}} = \bar x/s$.
\end{definition}

\begin{proposition}[Standard error of the Sharpe ratio]\label{prop:qm:estimation:srse}
For independent normal returns, $\sqrt n(\widehat{\mathrm{SR}} - \mathrm{SR}) \xrightarrow{d} \mathcal N(0, 1 + \tfrac12\mathrm{SR}^2)$ per period.
\end{proposition}

\begin{proof}
$(\bar x, s^2)$ are independent with asymptotic variances $\sigma^2$ and $2\sigma^4$. The gradient of $\mu/\sqrt v$ is $(1/\sigma,
-\mu/(2\sigma^3))$, so the \omterm{met:qm:estimation:delta}{delta method} gives $1 + \mu^2\cdot2\sigma^4/(4\sigma^6) = 1 + \tfrac12\mathrm{SR}^2$.
\end{proof}

Annualised, the \omterm{def:qm:estimation:estimator}{standard error} is about $\sqrt{(1 + \mathrm{SR}_1^2/2)/\text{years}}$, since the per-day \omterm{def:qm:estimation:sharpe}{Sharpe ratio} $\mathrm{SR}_1$ is small:
five years estimate an annual \omterm{def:qm:estimation:sharpe}{Sharpe ratio} to $\pm 0.45$, whatever its size. The researcher's strategy has $\widehat{\mathrm{SR}} = 1.51$
with that \omterm{def:qm:estimation:estimator}{standard error} under independence (Lo, 2002, derives the general case). With non-normal returns $B$ changes;
with autocorrelated returns, the annualisation by $\sqrt{252}$ and the variance both change, which is the next section's
subject: the HAC version of the same \omterm{met:qm:estimation:delta}{delta method} gives 0.87.

\section{Serial correlation: HAC variances}

The mean of a stationary series with autocovariances $\gamma(k)$ has
\[
\Var(\bar x) = \frac1n\sum_{|k|<n}\Bigl(1 - \frac{|k|}n\Bigr)\gamma(k) \approx \frac1n\sum_k\gamma(k) = \frac{\text{long-run variance}}n.
\]
The researcher's P\&L is the average of five overlapping five-day positions, so its autocorrelations are $1 - k/5$ for $k
\le 4$ (\cref{fig:qm:estimation:acf}), and the long-run variance is $1 + 2(0.8 + 0.6 + 0.4 + 0.2) = 5$ times the variance: the naive
\omterm{def:qm:estimation:estimator}{standard error} is $\sqrt5 = 2.24$ times too small.

\begin{definition}[HAC estimator]\label{def:qm:estimation:hac}
A heteroskedasticity- and autocorrelation-consistent (\emph{HAC estimator}\index{HAC estimator}) estimates a long-run
variance by weighting sample autocovariances; the Newey--West \omterm{def:qm:estimation:estimator}{estimator} $\hat\gamma(0) + 2\sum_{k=1}^L(1 - \frac k{L+1})\hat\gamma(k)$
uses Bartlett weights, which keep it nonnegative, with a bandwidth $L$ growing slowly with $n$ (a common rule is $L =
\lfloor4(n/100)^{2/9}\rfloor$).
\end{definition}

\begin{omfigure}
\begin{tikzpicture}
\begin{groupplot}[group style={group size=2 by 1,horizontal sep=1.6cm},
  width=7.6cm,height=4.8cm,tick label style={font=\small},label style={font=\small},
  grid=major,grid style={gray!25},table/col sep=comma]
\nextgroupplot[xlabel={lag (days)},ylabel={autocorrelation},xmin=0.3,xmax=10.7,ymin=-0.15,ymax=1,ybar,bar width=6pt,
  legend style={font=\footnotesize,at={(0.97,0.97)},anchor=north east,fill=white,draw=none}]
\addplot[fill=omDef!60,draw=omDef] table[x=lag,y=sample]{figdata/methods/11-estimation/acf.csv};
\addplot[fill=omThm!40,draw=omThm] table[x=lag,y=theory]{figdata/methods/11-estimation/acf.csv};
\legend{sample,$1 - k/5$}
\nextgroupplot[xlabel={Newey--West lags $L$},ylabel={se of the mean (bp)},xmin=0,xmax=30,ymin=0,ymax=5,ytick={0,1,2,3},
  legend style={font=\footnotesize,at={(0.97,0.97)},anchor=north east,fill=white,draw=none}]
\addplot[omDef,very thick] table[x=lags,y=nw]{figdata/methods/11-estimation/nwlags.csv};
\addplot[omThm,dashed,thick] table[x=lags,y=iid]{figdata/methods/11-estimation/nwlags.csv};
\addplot[black,dotted,thick] table[x=lags,y=true]{figdata/methods/11-estimation/nwlags.csv};
\legend{Newey--West,iid formula,true}
\end{groupplot}
\end{tikzpicture}

\omcaption{The researcher's five years of daily P\&L. Left: sample autocorrelations and the theoretical $1 - k/5$ of five
overlapping positions. Right: the Newey--West \omterm{def:qm:estimation:estimator}{standard error} of the mean P\&L against the number of lags: 1.19 bp with
none (the iid formula), 2.32 at the rule's seven lags, against a true 2.65. Data: the chapter's tutorial,
seeded.}\label{fig:qm:estimation:acf}
\end{omfigure}

With seven lags, Newey--West gives 2.32 basis points and a $t$-statistic of 1.73: the colleague's number. It still
understates the true 2.65, because the estimated autocovariances are biased toward zero in a sample of 1\,260 days. Over
2\,000 simulated histories of the strategy, nominal 95\% intervals built on the iid \omterm{def:qm:estimation:estimator}{standard error} contain the true mean
62.5\% of the time; with Newey--West, 92.2\% with seven lags and 93.5\% with twenty (\cref{fig:qm:estimation:coverage}). A HAC
\omterm{def:qm:estimation:estimator}{standard error} is necessary, not sufficient; the model-based alternative, when the overlap is known, uses the factor 5
directly.

\begin{omfigure}
\begin{tikzpicture}
\begin{axis}[width=9cm,height=4.2cm,xbar,bar width=10pt,
  xlabel={share of 95\% intervals containing the true mean (\%)},
  ytick={0,1,2},yticklabels={iid,{NW, 7 lags},{NW, 20 lags}},
  tick label style={font=\small},label style={font=\small},
  xmin=0,xmax=100,ymin=-0.6,ymax=2.6,grid=major,grid style={gray!25},
  table/col sep=comma]
\addplot[fill=omDef!60,draw=omDef,nodes near coords,nodes near coords style={font=\footnotesize,anchor=east,xshift=-1pt}] table[x=coverage,y=k]{figdata/methods/11-estimation/coverage.csv};
\addplot[black,dashed,thick,sharp plot] coordinates {(95,-0.6) (95,2.6)};
\end{axis}
\end{tikzpicture}

\omcaption{Coverage of nominal 95\% confidence intervals for the mean daily P\&L of the overlapping strategy, over 2\,000
simulated five-year histories: the iid \omterm{def:qm:estimation:estimator}{standard error} covers 62.5\% of the time, Newey--West 92--94\%. Data: the chapter's
tutorial, seeded.}\label{fig:qm:estimation:coverage}
\end{omfigure}

The same arithmetic tells a researcher how long to wait. An annual \omterm{def:qm:estimation:sharpe}{Sharpe ratio} $\mathrm{SR}$ estimated from independent daily
returns reaches a $t$-statistic of 2 after about $4/\mathrm{SR}^2$ years: one year at a \omterm{def:qm:estimation:sharpe}{Sharpe ratio} of 2, four at 1, sixteen at
0.5. If the long-run variance is $m$ times the variance, as with the overlapping positions ($m = 5$), the naive annualisation
overstates the \omterm{def:qm:estimation:sharpe}{Sharpe ratio} by $\sqrt m$ and the wait is $m$ times longer (\cref{fig:qm:estimation:years}): the researcher's
naive 1.51 needs 1.8 years, the honest reading 8.8. Five years of history cannot separate an honest \omterm{def:qm:estimation:sharpe}{Sharpe ratio} of
0.68 from zero.

\begin{omfigure}
\begin{tikzpicture}
\begin{semilogyaxis}[width=10cm,height=5cm,xlabel={naively annualised Sharpe ratio},ylabel={years to reach $t = 2$},
  tick label style={font=\small},label style={font=\small},xmin=0.5,xmax=3,ymin=0.3,ymax=100,
  ytick={0.5,1,2,5,10,20,50,100},log ticks with fixed point,grid=major,grid style={gray!25},
  legend columns=3,legend style={font=\footnotesize,at={(0.5,-0.32)},anchor=north,draw=none,fill=none,
  /tikz/every even column/.append style={column sep=8pt}},table/col sep=comma]
\addplot[omDef,very thick] table[x=sr,y=iid]{figdata/methods/11-estimation/years.csv};
\addplot[omThm,very thick,dashed] table[x=sr,y=overlap]{figdata/methods/11-estimation/years.csv};
\addplot[black,dotted,thick,sharp plot] coordinates {(0.5,5) (3,5)};
\addplot[only marks,mark=*,mark size=1.8pt,black] coordinates {(1.51,1.762) (1.51,8.811)};
\legend{independent returns,{five-day overlap, $m = 5$},five years}
\end{semilogyaxis}
\end{tikzpicture}

\omcaption{Years of daily data before a strategy's \omterm{def:qm:estimation:sharpe}{Sharpe ratio} reaches a $t$-statistic of 2, against the \omterm{def:qm:estimation:sharpe}{Sharpe ratio}
annualised by $\sqrt{252}$: with independent returns and with the researcher's five overlapping positions (long-run variance
five times the variance). Dots: the researcher's 1.51. Data: the chapter's tutorial module.}\label{fig:qm:estimation:years}
\end{omfigure}

\section{Tutorial: the $t$-statistic that halved}\label{tut:qm:estimation:hac}

\begin{tutorial}
\textbf{Goal.} Rebuild the researcher's and the colleague's \omterm{def:qm:estimation:estimator}{standard errors}, measure how often each is right, and see the
sandwich correct a misspecified likelihood. \textbf{End state:}
\cref{fig:qm:estimation:sandwich,fig:qm:estimation:acf,fig:qm:estimation:coverage} and the $t$-statistics 3.38 and 1.73.
\begin{tutsteps}
\item \textbf{The long-run variance} with Bartlett weights, and the \omterm{def:qm:estimation:estimator}{standard error} of a mean.
\omcode{code/firm/estim/firm_estim.py}{107}{124}{Newey--West long-run variance and the standard error of a mean.}\label{lst:qm:estimation:nw}
\item \textbf{The \omterm{def:qm:estimation:sharpe}{Sharpe ratio}} with its iid and HAC \omterm{def:qm:estimation:estimator}{standard errors}, by the \omterm{met:qm:estimation:delta}{delta method}.
\omcode{code/firm/estim/firm_estim.py}{127}{145}{The Sharpe ratio and its two standard errors.}\label{lst:qm:estimation:sharpe}
\item \textbf{Run} \texttt{problem()}, \texttt{coverage()}, \texttt{misspecified\_fit()} and \texttt{fig\_estimation.py}.
\end{tutsteps}
\textbf{What to change next.} Give the P\&L volatility clustering and compare the iid, White and Newey--West \omterm{def:qm:estimation:estimator}{standard errors};
estimate the mean from non-overlapping five-day returns and compare the efficiency.
\end{tutorial}

\section{Build: the estimation toolkit}\label{bld:qm:estimation:estim}

\begin{build}
\bfield{Purpose} Every estimate the miniature firm reports (a mean return, a \omterm{def:qm:estimation:sharpe}{Sharpe ratio}, a fitted parameter) carries a
\omterm{def:qm:estimation:estimator}{standard error} computed here, robust by default.
\bfield{Interface} \texttt{minimize(f, x0)}; \texttt{gradient}, \texttt{hessian}; \texttt{mle(loglik\_obs, x0)} returning the
estimate and its Hessian, outer-product and sandwich covariances; \texttt{long\_run\_variance(x, lags)}; \texttt{nw\_lags(n)};
\texttt{mean\_se(x, kind)}; \texttt{sharpe(x, periods)}; \texttt{delta\_method(g, theta, cov)}.
\bfield{Rules} Per-observation log-likelihoods (so scores can be formed); the sandwich is reported whenever a likelihood is
maximised; HAC bandwidth by the rule unless stated; no \omterm{def:qm:estimation:estimator}{standard error} without its method named in the output.
\bfield{Acceptance tests} \texttt{code/firm/estim/tests/}: the MLE of normal and exponential samples and their \omterm{def:qm:estimation:estimator}{standard errors};
the sandwich equals the Hessian variance for a correct model; Newey--West recovers the long-run variance of a moving average;
the Sharpe \omterm{def:qm:estimation:estimator}{standard error} matches simulation.
\bfield{Stretch} Andrews' automatic bandwidth; GMM with an efficient weight matrix and the overidentification test.
\end{build}

\begin{omsources}
\item P.~J.~Huber, ``The behavior of maximum likelihood estimates under nonstandard conditions'', \emph{Fifth Berkeley
Symposium}, 1967.
\item H.~White, ``Maximum likelihood estimation of misspecified models'', \emph{Econometrica} 50, 1982.
\item W.~K.~Newey and K.~D.~West, ``A simple, positive semi-definite, heteroskedasticity and autocorrelation consistent covariance
matrix'', \emph{Econometrica} 55, 1987.
\item L.~P.~Hansen, ``Large sample properties of generalized method of moments estimators'', \emph{Econometrica} 50, 1982.
\item A.~W.~Lo, ``The statistics of Sharpe ratios'', \emph{Financial Analysts Journal} 58, 2002.
\item G.~Casella and R.~L.~Berger, \emph{Statistical Inference}, Duxbury, 2nd ed., 2002.
\end{omsources}

\section{Exercises}

\begin{exercise}[$\star$]\label{exo:qm:estimation:1}
A hundred waiting times between trades sum to 50 seconds. Estimate the arrival rate and its \omterm{def:qm:estimation:estimator}{standard error}.
\end{exercise}

\begin{exercise}[$\star$]\label{exo:qm:estimation:2}
A strategy's annualised \omterm{def:qm:estimation:sharpe}{Sharpe ratio} is estimated at 1.0 from four years of independent daily returns. What is its \omterm{def:qm:estimation:estimator}{standard error}?
\end{exercise}

\begin{exercise}[$\star$]\label{exo:qm:estimation:3}
Is the sample variance with divisor $n$ unbiased? Consistent?
\end{exercise}

\begin{exercise}[$\star\star$]\label{exo:qm:estimation:4}
Daily P\&L follows an AR(1) with coefficient 0.2. By what factor does the naive \omterm{def:qm:estimation:estimator}{standard error} of the mean understate the
truth?
\end{exercise}

\begin{exercise}[$\star\star$]\label{exo:qm:estimation:5}
A signal was right on 110 of 200 trades. Using the \omterm{def:qm:estimation:mle}{Fisher information} of a Bernoulli variable, give the \omterm{def:qm:estimation:estimator}{standard error} of the hit
rate, and the $t$-statistic against 50\%.
\end{exercise}

\begin{exercise}[$\star\star$]\label{exo:qm:estimation:6}
For $n$ independent normal observations, use the \omterm{met:qm:estimation:delta}{delta method} to find the \omterm{def:qm:estimation:estimator}{standard error} of $\ln\hat\sigma$, and evaluate it at $n = 2\,000$.
\end{exercise}

\begin{exercise}[$\star\star\star$]\label{exo:qm:estimation:7}
\emph{Coding.} Over 2\,000 simulated histories of the overlapping strategy, measure the coverage of 95\% intervals built on the
iid, Newey--West(7) and Newey--West(20) \omterm{def:qm:estimation:estimator}{standard errors}.
\end{exercise}

\begin{exercise}[$\star\star\star$]\label{exo:qm:estimation:8}
\emph{Find the flaw.} ``Each day we record our strategy's trailing 20-day return. Over three years the \omterm{def:qm:estimation:sharpe}{Sharpe ratio} of these
daily figures, annualised by $\sqrt{252}$, is 1.4, with a \omterm{def:qm:estimation:estimator}{standard error} of $1/\sqrt3 = 0.58$: significant at the 5\% level.''
\end{exercise}

\section{Problem: The $t$-Statistic That Halved}

\begin{problem}[{Weekend problem --- overlapping positions and the standard error of a mean}]\label{pb:qm:estimation:1}
A strategy opens a position every day and holds it five days, so each day's P\&L is the average of five overlapping positions.
Over five years (1\,260 days) its daily P\&L, in basis points of capital, is the chapter's seeded history.

\medskip\noindent\textbf{Part I --- The naive view.}
\begin{enumerate}
\item What are the sample mean and standard deviation?
\item What are the iid \omterm{def:qm:estimation:estimator}{standard error} and $t$-statistic?
\item What are the annualised \omterm{def:qm:estimation:sharpe}{Sharpe ratio} and its iid \omterm{def:qm:estimation:estimator}{standard error}?
\item What are the first four sample autocorrelations, against theory?
\item Why are consecutive days correlated?
\end{enumerate}

\medskip\noindent\textbf{Part II --- Robust \omterm{def:qm:estimation:estimator}{standard errors}.}
\begin{enumerate}[resume]
\item What is the long-run variance factor of the P\&L, and the true \omterm{def:qm:estimation:estimator}{standard error} of the mean?
\item How many lags does the rule give, and what do Newey--West and the $t$-statistic become?
\item With twenty lags?
\item What is the HAC \omterm{def:qm:estimation:estimator}{standard error} of the \omterm{def:qm:estimation:sharpe}{Sharpe ratio}?
\item How often do iid 95\% intervals contain the true mean, over 2\,000 simulated histories?
\end{enumerate}

\medskip\noindent\textbf{Part III --- What the data can say.}
\begin{enumerate}[resume]
\item Knowing the overlap, what model-based \omterm{def:qm:estimation:estimator}{standard error} would you use?
\item What $t$-statistic should the researcher expect, on average, with a true mean of 4 bp?
\item How many years would a $t$-statistic of 3 need, on average?
\item Why is Newey--West still slightly too small?
\item What would evaluating non-overlapping five-day P\&L change?
\end{enumerate}

\medskip\noindent\textbf{Part IV --- Judgement.}
\begin{enumerate}[resume]
\item What other dependence in daily P\&L would the iid formula miss?
\item The researcher tried twenty variants before this one. What else must change in the evaluation?
\item What should the report on the strategy state?
\item State the \emph{named result}: the naive and the robust $t$-statistics.
\item In one sentence: when is the iid \omterm{def:qm:estimation:estimator}{standard error} of a mean wrong?
\end{enumerate}
\end{problem}

\section{Interview questions}

\begin{interviewq}[$\star$ \iqroles{researcher, mle}]\label{iq:qm:estimation:1}
What is maximum likelihood, and why is it the default \omterm{def:qm:estimation:estimator}{estimator}?
\end{interviewq}

\begin{interviewq}[$\star\star$ \iqroles{researcher, trader}]\label{iq:qm:estimation:2}
A strategy shows a \omterm{def:qm:estimation:sharpe}{Sharpe ratio} of 1.0 over one year. How confident are you that it is positive?
\end{interviewq}

\begin{interviewq}[$\star\star$ \iqroles{researcher}]\label{iq:qm:estimation:3}
What is a \omterm{def:qm:estimation:sandwich}{sandwich variance}, and when do you need it?
\end{interviewq}

\begin{interviewq}[$\star\star$ \iqroles{researcher}]\label{iq:qm:estimation:4}
You compute monthly returns every day from overlapping windows and regress them on a signal. What is wrong with the usual
\omterm{def:qm:estimation:estimator}{standard errors}, and how do you fix them?
\end{interviewq}

\begin{interviewq}[$\star\star$ \iqroles{researcher, mle}]\label{iq:qm:estimation:5}
What is the \omterm{def:qm:estimation:mle}{Fisher information}, and what does the Cramér--Rao bound say?
\end{interviewq}

\begin{interviewq}[$\star\star\star$ \iqroles{researcher}]\label{iq:qm:estimation:6}
If your likelihood is wrong, what does its maximiser estimate?
\end{interviewq}
